\documentclass[11pt]{article}
\usepackage[T1]{fontenc}
\usepackage{lmodern}
\usepackage{amsmath,amssymb,amsthm}
\usepackage[a4paper,margin=1in]{geometry}
\usepackage[hidelinks]{hyperref}
\usepackage{microtype}

\newtheorem{theorem}{Theorem}
\newtheorem{proposition}[theorem]{Proposition}
\newtheorem{corollary}[theorem]{Corollary}
\theoremstyle{definition}
\newtheorem{definition}[theorem]{Definition}
\theoremstyle{remark}
\newtheorem{remark}[theorem]{Remark}

\newcommand{\K}{\mathbb{K}}
\newcommand{\Q}{\mathbb{Q}}
\newcommand{\lam}{\lambda}
\newcommand{\etap}{\eta}
\newcommand{\taup}{\tau}
\newcommand{\Ha}{\widehat H}
\newcommand{\Da}{\widehat D}

\title{A Six-Parameter Extension of the\\
Direct Somos--4 Algebraic Generating Function}
\author{Thomas Scheuerle}
\date{}

\begin{document}
\maketitle

\begin{abstract}
Our earlier work constructed a five-parameter algebraic generating function
whose Hankel determinants start with $(\lambda,m,r,\eta)$ and obey a Somos--4
recurrence. The natural next step is to free the coefficient of the first
product term:
\[
 H_{k+4}H_k=\alpha H_{k+3}H_{k+1}+\tau H_{k+2}^2.
\]
This adds one more parameter, giving us six total. The construction itself is
straightforward: set $q^6=\alpha$ and rescale the previous five parameters by
suitable powers of $q$, then replace $x$ with $qx$. The result is a Hankel
sequence with the desired six parameters, while the first four determinants
remain exactly $\lambda, m, r, \eta$.

But here a subtle point arises. The determinant sequence itself is rational in
the six parameters, as is the quartic curve it generates. Yet our transported
generating function is not: the coefficient of $x$ becomes $q^{-5}r$, involving
a sixth root of $\alpha$. We can thus write the determinants and curve over the
rational field without adjoining any root, yet this particular power-series lift
requires one.

A comparison with Chang and Liu's quadratic-orthogonal-pair construction
illuminates these phenomena. With a natural shift, their bilateral Somos--4
sequence becomes ours, and their subdiagonal data satisfy $d_n=B_{n+2}$. Their
square-root choice $u^2=\alpha$ also admits a clean interpretation: replacing
$u$ by $-u$ corresponds exactly to replacing a generating series $F(x)$ by
$F(-x)$. This operation fixes every ordinary Hankel determinant but multiplies
the once-shifted determinant of order $N$ by $(-1)^N$. So ordinary Hankel
determinants lose a parity choice that the shifted ones preserve.
\end{abstract}

\section{Introduction}
The preceding paper~\cite{Scheuerle2026} started with four prescribed Hankel
determinants,
\[
 (H_1,H_2,H_3,H_4)=(\lambda,m,r,\eta),
\]
and constructed an algebraic generating function satisfying
\begin{equation}
 H_{k+4}H_k=H_{k+3}H_{k+1}+\tau H_{k+2}^2,
 \qquad k\geq1.
 \label{eq:old-somos4}
\end{equation}
Here we relax that construction by allowing the leading coefficient to vary:
\begin{equation}
 H_{k+4}H_k=\alpha H_{k+3}H_{k+1}+\tau H_{k+2}^2,
 \qquad k\geq1.
 \label{eq:new-somos4}
\end{equation}
We thus gain six independent parameters in place of five: $(\lambda,m,r,\eta,\alpha,\tau)$, with the special case $\alpha=1$ recovering the earlier construction.

Two questions present themselves, distinct in both scope and difficulty.

\emph{The existence question} is the simpler one: can we extend the old
construction so that the first four Hankel determinants remain exactly
$(\lambda,m,r,\eta)$ and the recurrence becomes \eqref{eq:new-somos4}?  The
answer is yes, and Section~2 provides an explicit weighted rescaling.

\emph{The rationality question} is harder. Can we write such a quadratic
algebraic generating function using only rational functions of the six
parameters over the field
\[
 K=\Q(\lambda,m,r,\eta,\alpha,\tau),
\]
while preserving the shifted-determinant identity $D_N=H_{N+2}$?  Our
weighted construction does not achieve this: it introduces a sixth root $q$ of
$\alpha$, and already its coefficient of $x$ is $q^{-5}r$. This is not a
cosmetic radical that simplifies away.

We need a few preliminary terms. By an \emph{ordinary Hankel sequence} we mean
the determinants $H_N$ of a coefficient sequence. A \emph{once-shifted
determinant} $D_N$ uses the same coefficients but with every index increased by
one. We call a coefficient sequence a \emph{moment sequence}, adopting the
language of continued fractions and orthogonal polynomials, though here
``moments'' are simply the coefficients of a formal power series—no measure is
involved. The \emph{Somos--4 orbit} is the sequence obtained by iterating the
recurrence for fixed initial values.

A key observation throughout this work is that three related objects carry
different information:
\[
 \text{Hankel determinant sequence},\qquad
 \text{algebraic curve},\qquad
 \text{generating series}.
\]
The first two are rational in the six parameters. The third need not be. We will
write ``descends to $K$'' as shorthand for ``can be written with coefficients in
$K$''. By this measure, the determinant sequence and the curve descend, while
our particular transported series does not.

This distinction clarifies our position relative to recent work. Liu, Wang, and
Zhang~\cite{LiuWangZhang2026} give broad rational-function constructions for
generic $(\alpha,\beta)$--Somos--4 Hankel transforms. Our problem is narrower:
we require four initial determinants to be independently prescribed, and we
insist on a specific shifted-determinant normalization. So the scope is more
restricted, though the geometric picture is richer.

Here is the plan. Section~2 performs the weighted rescaling and proves the
six-parameter recurrence. Section~3 tracks the shifted determinants, showing
exactly where fractional powers of $\alpha$ first appear. Section~4 establishes
that the determinant sequence and curve both descend to $K$. Section~5
translates our construction into Chang and Liu's notation and identifies the
Jacobi subdiagonal data. Section~6 interprets the sign choice as parity
conjugation. We conclude with a summary of the main facts.

\section{From the five-parameter construction to six parameters}
\subsection{The five-parameter input}
Let
\[
 G_{1;\lambda,m,r,\eta,\tau}(x)=\sum_{k\geq0}g_kx^k
\]
be the algebraic series from~\cite{Scheuerle2026}.  For any coefficient sequence
we define the ordinary and once-shifted Hankel determinants by
\[
 H_N(g)=\det(g_{i+j})_{0\leq i,j\leq N-1},\qquad
 D_N(g)=\det(g_{i+j+1})_{0\leq i,j\leq N-1}.
\]
That construction satisfies
\begin{equation}
 (H_1,H_2,H_3,H_4)=(\lambda,m,r,\eta),
 \qquad D_N=H_{N+2}.
 \label{eq:old-normalization}
\end{equation}
The underlying curve has discriminant
\begin{equation}
 \Delta_1(x)
 =\bigl(x\lambda^2\eta m r^2\bigr)^2
  \left((1-Jx+\tau x^2)^2-4x^3\right),
 \label{eq:old-discriminant}
\end{equation}
where
\begin{equation}
 J=\frac{\lambda\eta}{mr}
   +\frac{m^2}{\lambda r}
   +\frac{r^2}{m\eta}
   +\frac{\tau mr}{\lambda\eta}.
 \label{eq:old-J}
\end{equation}
For our purposes, these two facts suffice.

\subsection{Weighted transport}
The key idea is to introduce a sixth root of $\alpha$ and use it to rescale
carefully. Set
\[
 \K=\Q(\lambda,m,r,\eta,q,\tau),\qquad q^6=\alpha,
\]
and define new parameters
\begin{equation}
 \lambda'=\lambda,\qquad m'=q^{-2}m,\qquad
 r'=q^{-6}r,\qquad \eta'=q^{-12}\eta,\qquad
 \tau'=q^{-8}\tau.
 \label{eq:parameter-scaling}
\end{equation}
Now form the transported series
\begin{equation}
 G_{\alpha;\lambda,m,r,\eta,\tau}(x)
 :=G_{1;\lambda',m',r',\eta',\tau'}(qx).
 \label{eq:transported-G}
\end{equation}
If the original has expansion $G_{1;\lambda',m',r',\eta',\tau'}(x)=\sum_{k\geq0}a_kx^k$ and the transported has
$G_{\alpha;\lambda,m,r,\eta,\tau}(x)=\sum_{k\geq0}b_kx^k$, then
\begin{equation}
 b_k=q^ka_k.
 \label{eq:coefficient-scaling}
\end{equation}

\begin{proposition}[Determinant weights]
For every $N\geq1$,
\begin{equation}
 H_N(b)=q^{N(N-1)}H_N(a),\qquad
 D_N(b)=q^{N^2}D_N(a).
 \label{eq:determinant-weights}
\end{equation}
\end{proposition}
\begin{proof}
In the ordinary Hankel matrix, its $(i,j)$ entry is multiplied by $q^{i+j}$.
Factoring $q^i$ from row $i$ and $q^j$ from column $j$ gives the exponent
$2\sum_{j=0}^{N-1}j=N(N-1)$.  For the shifted matrix the entry is multiplied
by $q^{i+j+1}$.  The additional factor $q$ in every row contributes $N$, so
the total exponent is $N(N-1)+N=N^2$.
\end{proof}

\subsection{The six-parameter Somos--4 recurrence}
\begin{theorem}
Let $\Ha_N=H_N(b)$.  Then
\[
 (\Ha_1,\Ha_2,\Ha_3,\Ha_4)=(\lambda,m,r,\eta)
\]
and
\begin{equation}
 \Ha_{k+4}\Ha_k
 =\alpha\Ha_{k+3}\Ha_{k+1}+\tau\Ha_{k+2}^2,
 \qquad k\geq1.
 \label{eq:transported-recurrence}
\end{equation}
\end{theorem}
\begin{proof}
Equation \eqref{eq:determinant-weights} and
\eqref{eq:parameter-scaling} give
\[
 \Ha_1=\lambda,
 \quad \Ha_2=q^2m'=m,
 \quad \Ha_3=q^6r'=r,
 \quad \Ha_4=q^{12}\eta'=\eta.
\]
Write $H'_N=H_N(a)$.  The primed sequence satisfies
\[
 H'_{k+4}H'_k=H'_{k+3}H'_{k+1}+\tau'H'_{k+2}{}^2.
\]
The weight difference between the left-hand product and the first product on
the right is $6$, while that between the left-hand product and
$H'_{k+2}{}^2$ is $8$.  Hence
\[
 \Ha_{k+4}\Ha_k
 =q^6\Ha_{k+3}\Ha_{k+1}+q^8\tau'\Ha_{k+2}^2.
\]
Now use $q^6=\alpha$ and $q^8\tau'=\tau$.
\end{proof}

\begin{corollary}[Low-order recurrence check]
The first two values forced beyond the initial tuple are
\[
 \Ha_5=\frac{\alpha\eta m+\tau r^2}{\lambda},
\qquad
 \Ha_6=\frac{\alpha\Ha_5r+\tau\eta^2}{m}.
\]
These formulas provide a direct indexing check for any coefficient expansion
of \eqref{eq:transported-G}.
\end{corollary}

\section{Why fractional powers appear}
The weighted rescaling delivers the correct ordinary Hankel determinants, but
the shifted determinants behave differently. Understanding this gap is the key
to the rationality problem.

\subsection{The transported companion identity}
\begin{proposition}
If the primed series obeys $D_N(a)=H'_{N+2}$, then
\begin{equation}
 \Da_N:=D_N(b)=q^{-3N-2}\Ha_{N+2},
 \qquad N\geq1.
 \label{eq:graded-companion}
\end{equation}
\end{proposition}
\begin{proof}
By \eqref{eq:determinant-weights},
\[
 D_N(b)=q^{N^2}H'_{N+2},\qquad
 \Ha_{N+2}=q^{(N+2)(N+1)}H'_{N+2}.
\]
Their exponent difference is
$N^2-(N+2)(N+1)=-3N-2$.
\end{proof}

\begin{remark}
A constant coefficient rescaling $b_k\mapsto cb_k$ changes the ratio
$D_N(b)/\Ha_{N+2}$ by $c^{-2}$.  It cannot remove the factor $q^{-3N}$.
Thus the literal identity $D_N=H_{N+2}$ is not invariant under this graded
transport.  Equation \eqref{eq:graded-companion}, rather than the ungraded
identity, is the companion normalization naturally inherited from the
five-parameter series.
\end{remark}

\subsection{Parity conjugation and the companion orientation}
For any formal series
\[
 F(x)=\sum_{n\geq0}f_nx^n,
 \qquad
 P(F)(x):=F(-x)=\sum_{n\geq0}(-1)^nf_nx^n,
\]
the involution $P$ has different effects on the ordinary and once-shifted
Hankel determinants.

\begin{proposition}[Parity determinant laws]
For every $N\geq1$,
\begin{equation}
 H_N(P(F))=H_N(F),\qquad
 D_N(P(F))=(-1)^N D_N(F).
 \label{eq:parity-determinant-laws}
\end{equation}
\end{proposition}

\begin{proof}
The $(i,j)$ entry of the ordinary Hankel matrix is multiplied by
$(-1)^{i+j}$.  Factoring $(-1)^i$ from row $i$ and $(-1)^j$ from column $j$
gives the total exponent
$2\sum_{j=0}^{N-1}j=N(N-1)$, which is even.  In the once-shifted matrix the
entry is multiplied by $(-1)^{i+j+1}$.  The additional factor $-1$ in every
row contributes $(-1)^N$, proving the second identity.
\end{proof}

\begin{corollary}[Parity-conjugate companion law]
If $D_N(F)=c_NH_{N+2}(F)$, then
\begin{equation}
 D_N(P(F))=(-1)^Nc_NH_{N+2}(P(F)).
 \label{eq:parity-companion-law}
\end{equation}
In particular, the ordinary Hankel orbit cannot distinguish $F(x)$ from
$F(-x)$, whereas a prescribed companion law distinguishes the two parity
orientations whenever some odd-order companion determinant is nonzero.
\end{corollary}

Applied to our transported series, parity conjugation preserves the
six-parameter Somos--4 orbit yet alters the graded companion identity:
\begin{equation}
 D_N(G_\alpha(-x))=(-1)^Nq^{-3N-2}H_{N+2}(G_\alpha(-x)).
 \label{eq:transported-parity-companion}
\end{equation}
The upshot is concrete: $F(x)$ and $F(-x)$ share the same ordinary Hankel
determinants but differ by a predictable sign in their shifted determinants.
We reserve the term ``orientation'' for this sign choice alone. No interpretation
via algebraic sheets or bilateral Somos reversal is required.

\subsection{The first coefficients and failure of descent}
The companion normalization of the five-parameter series gives $g_0=\lambda$ and
$g_1=r$. When we apply the rescaling \eqref{eq:parameter-scaling}, we get
$a_0=\lambda$ and $a_1=r'=q^{-6}r$. By \eqref{eq:coefficient-scaling},
\begin{equation}
 b_0=\lambda,\qquad b_1=q^{-5}r=\frac{qr}{\alpha}.
 \label{eq:first-coefficients}
\end{equation}
Expanding further:
\begin{align}
 b_2&=\frac{m}{\lambda}+\frac{q^{-10}r^2}{\lambda}
     =\frac{m}{\lambda}+\frac{q^2r^2}{\alpha^2\lambda},
 \label{eq:b2}\\
 b_3&=\frac{q^5m^2}{\lambda^2r}
 +\frac{q^3\eta}{\alpha r}
 +\frac{2qmr}{\alpha\lambda^2}
 +\frac{q^3r^3}{\alpha^3\lambda^2}.
 \label{eq:b3}
\end{align}
Putting these together:
\begin{equation}
\begin{aligned}
G_\alpha(x)={}&\lambda+\frac{qr}{\alpha}x
 +\left(\frac{m}{\lambda}+\frac{q^2r^2}{\alpha^2\lambda}\right)x^2\\
&+\left(\frac{q^5m^2}{\lambda^2r}
 +\frac{q^3\eta}{\alpha r}
 +\frac{2qmr}{\alpha\lambda^2}
 +\frac{q^3r^3}{\alpha^3\lambda^2}\right)x^3+O(x^4).
\end{aligned}
\label{eq:series-expansion}
\end{equation}
The coefficient $b_1$ also emerges from \eqref{eq:graded-companion} at $N=1$:
\[
 b_1=D_1(b)=q^{-5}H_3(b)=q^{-5}r.
\]
So the fractional power is not an artifact of poor simplification; it is forced
directly by the shifted-determinant relation at the first step.

Put
\[
 K=\Q(\lambda,m,r,\eta,\alpha,\tau),\qquad
 L=K(q),\qquad q^6=\alpha.
\]
\begin{proposition}[Non-descent of the transported branch]
In the generic case $r\neq0$, the transported series \eqref{eq:transported-G}
does not belong to $K[[x]]$.
\end{proposition}
\begin{proof}
The coefficient $b_1=q^{-5}r=qr/\alpha$ lies outside $K$ generically. By
Eisenstein's criterion, the polynomial $X^6-\alpha$ is irreducible over $K$, so
$q\notin K$. Since $r\neq0$, we have $b_1\notin K$, and thus the series does not
lie in $K[[x]]$.
\end{proof}

The radicals do not descend to $K$, yet they vanish from the ordinary Hankel
determinants. For instance,
\[
 H_2(b)=b_0b_2-b_1^2
 =\lambda\left(\frac{m}{\lambda}+\frac{q^{-10}r^2}{\lambda}\right)-q^{-10}r^2=m.
\]
This example crystallizes the distinction: fractional powers appear in the
coefficients but cancel in the determinants.

\section{What is defined over the six-parameter field}
The non-descent we proved applies to one particular generating series. It says
nothing about the recurrence itself or its algebraic curve. Let us check each
separately.

\subsection{Transport of the discriminant}
To track how the discriminant transforms, decompose the invariant $J$ as
\begin{equation}
 A=\frac{\lambda\eta}{mr}+\frac{\tau mr}{\lambda\eta},
 \qquad
 B=\frac{m^2}{\lambda r}+\frac{r^2}{m\eta}.
 \label{eq:A-B}
\end{equation}
Under the rescaling \eqref{eq:parameter-scaling},
\begin{equation}
 J'=q^{-4}A+q^2B.
 \label{eq:J-prime}
\end{equation}
The substitution $x\mapsto qx$ then gives
\begin{equation}
 qJ'=q^{-3}A+q^3B
 =\alpha^{-1/2}A+\alpha^{1/2}B.
 \label{eq:J-alpha}
\end{equation}

\begin{proposition}[Transported discriminant]
The discriminant obtained by substituting
\eqref{eq:parameter-scaling} and $x\mapsto qx$ into
\eqref{eq:old-discriminant} is
\begin{equation}
 \Delta_\alpha(x)
 =q^{-50}\bigl(x\lambda^2\eta m r^2\bigr)^2
 \left[
  \left(1-(q^{-3}A+q^3B)x+q^{-6}\tau x^2\right)^2
  -4q^3x^3
 \right].
 \label{eq:q-discriminant}
\end{equation}
Equivalently,
\begin{equation}
 \Delta_\alpha(x)
 =\alpha^{-25/3}\bigl(x\lambda^2\eta m r^2\bigr)^2
 \left[
  \left(1-(\alpha^{-1/2}A+\alpha^{1/2}B)x
             +\frac{\tau}{\alpha}x^2\right)^2
  -4\alpha^{1/2}x^3
 \right].
 \label{eq:alpha-discriminant}
\end{equation}
\end{proposition}
\begin{proof}
The square prefactor has total $q$-weight
$2(1-12-2-12)=-50$.  The three terms in the residual factor transform as
$J'x\mapsto qJ'x$, $\tau'x^2\mapsto q^2\tau'x^2=q^{-6}\tau x^2$, and
$x^3\mapsto q^3x^3$.  Equations \eqref{eq:J-prime} and $q^6=\alpha$ give the
stated forms.
\end{proof}

The overall monomial factor in \eqref{eq:q-discriminant} is a square in
$\K(x)$ and does not affect the residual double cover.  It should, however,
be retained when comparing discriminants of particular polynomial
representatives, because multiplying a quadratic equation by a scalar squares
that scalar in its discriminant.

\subsection{A rational model of the curve}
Dropping the square prefactor from \eqref{eq:alpha-discriminant} and setting
\begin{equation}
 x=\alpha^{1/2}z,
 \label{eq:x-z}
\end{equation}
the residual part becomes
\begin{equation}
 Y^2=\left(1-(A+\alpha B)z+\tau z^2\right)^2
      -4\alpha^2z^3.
 \label{eq:rational-quartic-z}
\end{equation}
Every coefficient here lies in $\Q(\lambda,m,r,\eta,\alpha,\tau)$.

To match the quartic convention from our earlier work~\cite{Scheuerle2026},
let $X=z^{-1}$ and $y=z^{-2}Y$, or equivalently $Y=z^2y$. Then
\eqref{eq:rational-quartic-z} is birational to
\begin{equation}
 y^2=\left(X^2-(A+\alpha B)X+\tau\right)^2-4\alpha^2X.
 \label{eq:rational-quartic-X}
\end{equation}
Thus the curve descends to the six-parameter field. However, the change of
variables \eqref{eq:x-z} rescales power-series coefficients and destroys the
Hankel normalization. So removing radicals from the curve is not the same as
removing them from a specific generating series.

\subsection{Formal representatives with the same Hankel determinants}
Hankel determinants do not uniquely specify a generating series. For a
nondegenerate sequence, define
\begin{equation}
 \beta_n=\frac{H_{n-1}H_{n+1}}{H_n^2},\qquad H_0=1.
 \label{eq:jacobi-beta}
\end{equation}
All $H_n$ from our recurrence, and hence all $\beta_n$, lie in $K$. For any
choice of $a_0,a_1,\ldots\in K$, the Jacobi continued fraction
\begin{equation}
 \widetilde G(x)=
 \cfrac{H_1}{1-a_0x-
  \cfrac{\beta_1x^2}{1-a_1x-
   \cfrac{\beta_2x^2}{1-a_2x-\ddots}}}
 \label{eq:jacobi-representative}
\end{equation}
has Hankel determinants $H_n$. So multiple formal power series over $K$
realize the same Hankel sequence. What they may lack is the special properties
we care about: being quadratic algebraic, satisfying the graded companion law,
or being the distinguished Stieltjes choice.

\section{Comparison with the construction of Chang and Liu}
We have shown that our transported series involves fractional powers of $\alpha$,
at least in the particular form we constructed it. A different quadratic series
might conceivably avoid them. Chang and Liu provide a natural comparison, since
their construction realizes generic bilateral Somos--4 sequences via Hankel
determinants. We now express both constructions in a common notation.

\subsection{The shifted Somos--4 orbit}
Chang and Liu \cite[Sections~3.3--4.2]{ChangLiu2026} treat the generic bilateral Somos--4 initial-value problem by Hankel determinants of quadratic orthogonal pairs. Their normalized recurrence is
\begin{equation}
 \sigma_{n+2}\sigma_{n-2}=a\sigma_{n+1}\sigma_{n-1}+b\sigma_n^2,
 \qquad n\in\mathbb Z,
 \label{eq:CL-recurrence}
\end{equation}
with $\sigma_{-1}=x$, $\sigma_0=1$, $\sigma_1=y$, and $\sigma_2=z$.
Our freely prescribed determinants fit this normalization after the shift
\begin{equation}
 \sigma_n=\frac{H_{n+2}}m,
 \qquad
 (x,y,z;a,b)=\left(\frac\lambda m,\frac r m,\frac\eta m;\alpha,\tau\right).
 \label{eq:CL-dictionary}
\end{equation}
This shift uses all six parameters.  In particular, it avoids the unwanted constraint
$\eta=\alpha\lambda r+\tau m^2$ that would follow from imposing our recurrence at $N=0$ with $H_0=1$.

Write $u^2=\alpha$ and define
\begin{align}
 \mathcal N&=\tau r^2m+\lambda^2\eta^2+\alpha\eta m^3+\alpha\lambda r^3,
 \label{eq:CL-Nplus}\\
 \mathcal N_0&=\tau r^2m-\lambda^2\eta^2+\alpha\eta m^3+\alpha\lambda r^3.
 \label{eq:CL-Nminus}
\end{align}
Under \eqref{eq:CL-dictionary}, the Chang--Liu parameters include
\begin{align}
 d_0&=\frac{\lambda r}{m^2},&u_0&=\frac{2r}{m},\\
 \frac vu&=\frac{\mathcal N}{2\alpha\lambda r\eta m},&
 uv&=\frac{\mathcal N}{2\lambda r\eta m},\\
 \frac{v_0}{u}&=\frac{\mathcal N_0}{2\alpha\lambda r\eta m},&
 uv_0&=\frac{\mathcal N_0}{2\lambda r\eta m},
 \label{eq:CL-parameters}
\end{align}
and
\begin{equation}
 f=\frac{\tau}{\alpha}-
 \frac{\mathcal N^2}{4\alpha\lambda^2r^2\eta^2m^2}.
 \label{eq:CL-f}
\end{equation}
The sign difference between $\mathcal N$ and $\mathcal N_0$ is essential because $v_0$ enters the distinguished moment branch.
The scaling $X=uZ$ turns their curve into
\begin{equation}
 Y^2=(\alpha Z^2+f)^2+4\alpha Z-4uv,
 \label{eq:CL-descended-curve}
\end{equation}
which is defined over $K$. Again, this statement concerns the curve only.  It says nothing yet about whether the coefficients of a chosen generating series lie in $K$.

\subsection{Translation to Jacobi data}
Use the local Jacobi convention
\begin{equation}
 F(t)=\cfrac{\lambda}{1-A_0t-
 \cfrac{B_1t^2}{1-A_1t-
 \cfrac{B_2t^2}{1-A_2t-
 \cfrac{B_3t^2}{1-A_3t-B_4t^2M_4(t)}}}},
 \qquad M_4(0)=1.
 \label{eq:local-jacobi-convention}
\end{equation}
The product formula
\begin{equation}
 H_N=\lambda^N\prod_{j=1}^{N-1}B_j^{N-j}
 \label{eq:J-Hankel-product}
\end{equation}
gives
\begin{align}
 B_1&=\frac{m}{\lambda^2},&
 B_2&=\frac{\lambda r}{m^2},&
 B_3&=\frac{m\eta}{r^2},&
 B_4&=\frac{r(\alpha\eta m+\tau r^2)}{\lambda\eta^2}.
 \label{eq:first-B}
\end{align}
The companion identities through $D_3=H_5$ force
\begin{align}
 A_0&=\frac r\lambda,&
 A_1&=\frac{\eta\lambda^2+m^2}{\lambda mr},&
 A_2&=\frac{\alpha\eta m^3+\lambda r^3+\tau m^2r^2}{\eta\lambda mr}.
 \label{eq:first-A}
\end{align}

\begin{proposition}[Fourth shifted determinant]
For the convention \eqref{eq:local-jacobi-convention},
\begin{equation}
 D_4=\lambda^4B_1^3B_2^2B_3
 \left(A_0A_1A_2A_3-A_0A_1B_3-A_0A_3B_2-A_2A_3B_1+B_1B_3\right).
 \label{eq:generic-D4}
\end{equation}
After substitution of \eqref{eq:first-B} and \eqref{eq:first-A}, this reduces to
\begin{equation}
 D_4=\frac{\eta}{\lambda r^2}
 \left[A_3r(\alpha\eta m+\tau r^2)-\eta^2\lambda\right].
 \label{eq:D4-reduced}
\end{equation}
Consequently $D_4=H_6$, where
\[
 H_6=\frac{\alpha^2\eta mr+\alpha\tau r^3+\tau\eta^2\lambda}{\lambda m},
\]
uniquely forces
\begin{equation}
 A_3=
 \frac{\alpha^2\eta mr^3+\alpha\tau r^5+\eta^3\lambda m+\tau\eta^2\lambda r^2}
 {\eta mr(\alpha\eta m+\tau r^2)}.
 \label{eq:A3}
\end{equation}
\end{proposition}
\begin{proof}
Equation \eqref{eq:generic-D4} follows by direct expansion of the shifted $4\times4$ Hankel determinant in Jacobi parameters.  Substitution and collection of the terms linear in $A_3$ gives \eqref{eq:D4-reduced}.  Solving \eqref{eq:D4-reduced} for $A_3$ after inserting the displayed value of $H_6$ gives \eqref{eq:A3}.
\end{proof}
An equivalent useful form is
\[
 A_3=\frac{\alpha r^2}{\eta m}+
 \frac{\eta\lambda(\eta m+\tau r^2)}{mr(\alpha\eta m+\tau r^2)}.
\]
No $B_4$ occurs in $D_4$: the fourth ordinary Hankel step determines $B_4$, while the fourth shifted step determines $A_3$, without imposing any extra relation among the six parameters.

Chang and Liu use
\[
 d_n=\frac{\sigma_{n-1}\sigma_{n+1}}{\sigma_n^2}.
\]
Substitution of \eqref{eq:CL-dictionary} gives the exact all-order interface
\begin{equation}
 d_n=\frac{H_{n+1}H_{n+3}}{H_{n+2}^2}=B_{n+2}.
 \label{eq:d-B}
\end{equation}
Thus $d_0=B_2$, $d_1=B_3$, and $d_2=B_4$.  Thus the subdiagonal parameters match after an index shift: only $B_1$ comes before the Chang--Liu sequence begins.  The diagonal parameters and the normalization of the generating-series coefficients still have to be treated separately.

\section{The parity involution in the Chang--Liu branch}
The same parity symmetry appears in Chang and Liu's construction, though dressed
as a choice of square root. Let
\[
 S(t;u)=\sum_{n\geq0}s_n(u)t^n,\qquad u^2=\alpha,
\]
be their distinguished moment branch. Chang and Liu show (Remark~4.7):
\begin{equation}
 S(t;-u)=S(-t;u),\qquad s_n(-u)=(-1)^ns_n(u).
 \label{eq:galois-parity}
\end{equation}
The involution $u\mapsto-u$ is precisely the moment-parity involution from
Proposition~\ref{eq:parity-determinant-laws}. Consequently,
\begin{equation}
 H_N(S(t;-u))=H_N(S(t;u)),\qquad
 D_N(S(t;-u))=(-1)^ND_N(S(t;u)).
 \label{eq:CL-parity-determinants}
\end{equation}
So the two sign choices for $u$ do not yield two different ordinary Somos--4
sequences. Instead, they give coefficient sequences related by alternating signs.
Ordinary Hankel determinants cannot see this difference, but shifted ones can.

To see this directly in the coefficients, write
\[
 S(t;u)=E(t^2)+utR(t^2),
\]
where $E$ and $R$ have even-indexed moment series. By \eqref{eq:galois-parity},
the even moments are invariant under $u\mapsto-u$ while the odd moments change
sign. So the appearance of $u=\sqrt\alpha$ in this moment branch records the
parity orientation of the odd moments. This observation applies to the
particular lift we have chosen; it does not claim that every representation of
the ordinary Hankel orbit requires a quadratic extension.

The shifted-determinant normalization makes the choice concrete. If one branch
satisfies $D_N=c_NH_{N+2}$, its parity conjugate satisfies
$D_N=(-1)^Nc_NH_{N+2}$. Picking a literal or graded companion law thus selects
a parity orientation. This is all the involution contributes here.

\section{Summary}
Let us collect the main threads.

Our weighted transport \eqref{eq:transported-G} yields a quadratic algebraic
series over $K(q)$ (with $q^6=\alpha$) whose first four ordinary Hankel
determinants are $(\lambda,m,r,\eta)$ and whose full Hankel orbit satisfies
\eqref{eq:transported-recurrence}.

The once-shifted determinants obey the graded companion law
\eqref{eq:graded-companion}. Already the coefficient $b_1=q^{-5}r$ lies outside
$K$ generically, proving that this particular transported branch does not
descend.

Yet the curve itself descends: it has the rational model \eqref{eq:rational-quartic-z}
over $K$. This shows that the orbit and the curve can descend even when the
generating series does not.

The dictionary \eqref{eq:CL-dictionary} links our orbit to Chang and Liu's
construction, with subdiagonal interface $d_n=B_{n+2}$ holding throughout.
The involution $u\mapsto-u$ in their moment branch is precisely moment parity:
it fixes every ordinary Hankel determinant but multiplies each once-shifted
determinant of order $N$ by $(-1)^N$.

\section{Conclusion}
The construction is elegant and the limitation is real.

On the positive side, the weighted substitution \eqref{eq:transported-G} with
$q^6=\alpha$ extends our five-parameter generating function to the full
six-parameter Somos--4 recurrence. The first four determinants stay at
$(\lambda,m,r,\eta)$, and the whole sequence satisfies
\eqref{eq:transported-recurrence}.

On the other hand, generating-series coefficients carry more information than
Hankel determinants. Under the same rescaling, the shifted determinants follow
\eqref{eq:graded-companion}, and immediately $b_1=q^{-5}r$ escapes from
$K=\Q(\lambda,m,r,\eta,\alpha,\tau)$ in the generic case. The Hankel sequence
itself lies in $K$, and the quartic \eqref{eq:rational-quartic-z} shows the
curve does as well. Three seemingly parallel rationality questions have
different answers: the determinants and curve are rational in the six parameters,
but this particular power series is not.

The Chang--Liu comparison clarifies the remaining ambiguity. Their identity
$S(t;-u)=S(-t;u)$ means that flipping the sign of $u=\sqrt\alpha$ merely
alternates the signs of odd coefficients. This preserves all ordinary Hankel
determinants but changes each once-shifted determinant of order $N$ by
$(-1)^N$. In short: ordinary Hankel determinants lose a parity choice that the
shifted determinants retain. No deeper geometric interpretation is needed.

\section*{Disclosure on authoring and verification tools}
The manuscript was revised with assistance from AI-based language and symbolic
calculation tools.  Their use was directed toward notation control, LaTeX
editing, and exact algebraic checks.  The author remains responsible for the
mathematical statements and the final text.

\begin{thebibliography}{99}
\bibitem{ChangLiu2026}
X.-K.~Chang and J.~Liu,
\newblock Hankel determinants from quadratic orthogonal pairs for hyperelliptic functions and their applications,
\newblock arXiv:2603.11670v4 [nlin.SI], 2026; to appear in the \emph{Journal of the London Mathematical Society}.
\bibitem{LiuWangZhang2026}
F.~Liu, Y.~Wang, and Z.~Zhang,
\newblock Hankel transform and $(\alpha,\beta)$ Somos--4 sequences,
\newblock arXiv:2608.08703v1 [math.CO], 2026.
\bibitem{Scheuerle2026}
T.~Scheuerle,
\newblock Direct generation of a Somos--4 sequence from an algebraic generating function,
\newblock arXiv:2609.15754 [math.CO], 2026.

\bibitem{Wall1948}
H.~S. Wall,
\newblock \emph{Analytic Theory of Continued Fractions},
\newblock Van Nostrand, New York, 1948.

\bibitem{Flajolet1980}
P.~Flajolet,
\newblock Combinatorial aspects of continued fractions,
\newblock \emph{Discrete Mathematics} \textbf{32} (1980), 125--161.

\bibitem{Layman2001}
J.~W. Layman,
\newblock The Hankel transform and some of its properties,
\newblock \emph{Journal of Integer Sequences} \textbf{4} (2001), Article 01.1.5.

\bibitem{French2007}
C.~French,
\newblock Transformations preserving the Hankel transform,
\newblock \emph{Journal of Integer Sequences} \textbf{10} (2007), Article 07.7.3.

\bibitem{Hone2005}
A.~N.~W. Hone,
\newblock Elliptic curves and quadratic recurrence sequences,
\newblock \emph{Bulletin of the London Mathematical Society} \textbf{37} (2005), 161--171.

\bibitem{FominZelevinsky2002}
S.~Fomin and A.~Zelevinsky,
\newblock The Laurent phenomenon,
\newblock \emph{Advances in Applied Mathematics} \textbf{28} (2002), 119--144.

\bibitem{VanderPoortenSwart2006}
A.~J. van der Poorten and C.~S. Swart,
\newblock Recurrence relations for elliptic sequences: every Somos--4 is a Somos--$k$,
\newblock \emph{Bulletin of the London Mathematical Society} \textbf{38} (2006), 546--554.
\end{thebibliography}
\end{document}
